\subsection{APPROXIMATIONS TO A REAL NUMBER}

DEFINITION 1. Given a number \(\varepsilon > 0\) , a real number \(r\) is said to be an approximation to within \(\varepsilon\) to a real number \(x\) , if \(|x - r| \leqslant \varepsilon\) ; \(r\) is said to be an approximation by defect if \(r \leqslant x\) , by excess if \(r \geqslant x\) .

Let A be a dense subset of R. For each \(x \in R\) and each \(\varepsilon > 0\) there is an approximation by defect (resp. excess) to x to within \(\varepsilon\) belonging to A, since the interval \(]x - \varepsilon, x[\) (resp. \(]x, x + \varepsilon[\) ) contains at least one point of A. If we now consider a given strictly decreasing sequence ( \(\varepsilon_{n}\) ) of numbers \textgreater{} 0, tending to 0, and if \(r_{n}\) is an approximation to x to within \(\varepsilon_{n}\) , then the sequence ( \(r_{n}\) ) has x as limit as n tends to infinity.

In the case where A is a subgroup of the additive group R, and we restrict the \(\varepsilon_{n}\) to belong to A, we can define canonically for each \(x \in R\) a sequence \((r_{n})\) of approximations to x by defect, belonging to A.

For by Archimedes' axiom (§ 2, no. 1, Theorem 1), the set of integers \(p\) such that \(p\varepsilon_n \leqslant x\) has a greatest element \(p_n\) ; in other words, there is a unique integer \(p_n\) such that

\[
p _ {n} \varepsilon_ {n} \leqslant x <   (p _ {n} + 1) \varepsilon_ {n}.\tag{1}
\]

Since \(|x - p_n\varepsilon_n| \leqslant \varepsilon_n\) , it follows that \(p_n\varepsilon_n\) is an approximation to \(x\) to within \(\varepsilon_n\) by defect, and belongs to A by hypothesis; similarly \((p_n + 1)\varepsilon_n\) is an approximation to \(x\) to within \(\varepsilon_n\) by excess, belonging to A, and the two sequences \((p_n\varepsilon_n)\) and \(((p_n + 1)\varepsilon_n)\) have \(x\) as their limit.

